%! Parametrically Defined Circles and Lines — Unit 4, Lesson 4.4 — Group Activity
\documentclass[10pt]{article}
\usepackage{precalculus-article}
\usepackage{precalculus-boxes}
\newcommand{\TallMath}[1]{$\displaystyle #1\rule[-1.4em]{0pt}{3.2em}$}

\begin{document}

\pageheader{Unit 4, Lesson 4.4}{Group Activity}
\namepartnerperiod

\vspace{0.10in}

%----------------------------------------------------------------------
% Part A — Circular Motion
%----------------------------------------------------------------------
\begin{scenariobox}[Part A --- Satellite Orbit]{sky}
A satellite orbits Earth in a circular orbit of radius 8000 km, centered
at Earth's center (the origin). It starts at $(8000,\,0)$ and travels
counterclockwise, completing one full orbit in 2 hours. Using $t$ in
hours, the angle after $t$ hours is $\pi t$ radians.
\end{scenariobox}

\vspace{0.08in}

\begin{tcolorbox}[colback=white, colframe=black!40,
    title=\textbf{Tier R --- Remediate},
    fonttitle=\bfseries, arc=2mm, left=3mm, right=3mm, top=2mm, bottom=2mm]

\textbf{A1.}\quad Write parametric equations $x(t)$ and $y(t)$ for
the satellite's position at time $t$ hours.

\vspace{4pt}
$x(t) = $ \blank{5cm} \qquad $y(t) = $ \blank{5cm}

\vspace{0.12in}
\textbf{A2.}\quad Complete the table. (Round to the nearest whole number
where needed.)

\vspace{4pt}
\begin{center}
\renewcommand{\arraystretch}{1.7}
\begin{tabular}{|c|c|c|c|}
\hline
$t$ (hours) & $x(t)$ (km) & $y(t)$ (km) & Position $(x,\;y)$ \\
\hline
$0$   & \blank{1.8cm} & \blank{1.8cm} & \blank{2.8cm} \\
\hline
$0.5$ & \blank{1.8cm} & \blank{1.8cm} & \blank{2.8cm} \\
\hline
$1$   & \blank{1.8cm} & \blank{1.8cm} & \blank{2.8cm} \\
\hline
$1.5$ & \blank{1.8cm} & \blank{1.8cm} & \blank{2.8cm} \\
\hline
$2$   & \blank{1.8cm} & \blank{1.8cm} & \blank{2.8cm} \\
\hline
\end{tabular}
\end{center}

\end{tcolorbox}

\vspace{0.08in}

\begin{tcolorbox}[colback=white, colframe=black!40,
    title=\textbf{Tier A --- Approaching Proficiency},
    fonttitle=\bfseries, arc=2mm, left=3mm, right=3mm, top=2mm, bottom=2mm]

\textbf{A3.}\quad At what time $t$ (in hours) is the satellite directly
above the origin on the north side (i.e., $x = 0$ and $y > 0$)?
Show your reasoning.
\writelines{3}

\vspace{0.04in}
\textbf{A4.}\quad Write new parametric equations for a \emph{clockwise}
orbit of the same satellite along the same circular path.
\writelines{2}

\end{tcolorbox}

\vspace{0.08in}

%----------------------------------------------------------------------
% Part B — Linear Motion
%----------------------------------------------------------------------
\begin{scenariobox}[Part B --- Rover Path]{sky}
A rover travels in a straight line from base camp at $(-3,\,2)$ to a
sample site at $(5,\,-2)$.
\end{scenariobox}

\vspace{0.08in}

\begin{tcolorbox}[colback=white, colframe=black!40,
    title=\textbf{Tier E --- Extension},
    fonttitle=\bfseries, arc=2mm, left=3mm, right=3mm, top=2mm, bottom=2mm]

\textbf{B1.}\quad Write parametric equations for the rover's position
for $0 \le t \le 1$.
\writelines{2}

\vspace{0.04in}
\textbf{B2.}\quad Find the rover's position at $t = 0.25$ and at $t = 0.75$.
\writelines{3}

\vspace{0.04in}
\textbf{B3.}\quad Determine the value of $t$ when the rover crosses the
$x$-axis (where $y = 0$). Show your algebra, then state the rover's
coordinates at that moment.
\writelines{4}

\end{tcolorbox}

\end{document}
